\documentclass[11pt,a4paper]{article}

\usepackage[utf8]{inputenc}
\usepackage[indonesian]{babel}
\usepackage[a4paper,top=2cm,bottom=2cm,left=2cm,right=2cm]{geometry}
\usepackage{amsmath,amsfonts,amssymb}
\usepackage{booktabs}
\usepackage{enumitem}
\usepackage{xcolor}
\usepackage{colortbl}
\usepackage{tcolorbox}
\usepackage{fancyhdr}
\usepackage{titlesec}
\usepackage[breaklinks=true,colorlinks=true,linkcolor=primaryblue,citecolor=primaryblue,urlcolor=primaryblue]{hyperref}

\definecolor{primaryblue}{RGB}{24, 43, 73}
\definecolor{secondaryblue}{RGB}{41, 128, 185}
\definecolor{boxbg}{RGB}{245, 248, 252}

\pagestyle{fancy}
\fancyhf{}
\fancyhead[L]{\small\textbf{\color{primaryblue}Problem Set Minggu 9: Analisis Spektral \& Dekomposisi Wavelet}}
\fancyhead[R]{\small\color{gray}TEE-412 / Teknik Elektro FT UNIB}
\fancyfoot[L]{\footnotesize\color{gray}Dosen: Ir. Novalio Daratha, S.T., M.Sc., Ph.D.}
\fancyfoot[R]{\footnotesize Halaman \thepage\ dari 3}
\renewcommand{\headrulewidth}{0.5pt}
\renewcommand{\footrulewidth}{0.4pt}

\titleformat{\section}{\color{primaryblue}\normalfont\large\bfseries}{\thesection.}{0.5em}{}[\titlerule]
\setlist{nosep}

\begin{document}

\begin{center}
    {\large \textbf{\color{primaryblue}UNIVERSITAS BENGKULU -- FAKULTAS TEKNIK}}\\[0.1cm]
    {\large \textbf{\color{secondaryblue}JURUSAN TEKNIK ELEKTRO}}\\[0.2cm]
    {\LARGE \textbf{\color{primaryblue}PROBLEM SET 9: ANALISIS DERET WAKTU DAN PERAMALAN}}\\[0.1cm]
    {\normalsize \textbf{Topik: Transformasi Fourier (FFT), Spektrum Daya (PSD), dan Wavelet Diskrit (DWT)}}
\end{center}

\vspace{0.1cm}

\begin{tcolorbox}[colback=boxbg,colframe=primaryblue,arc=2mm,title=\textbf{Informasi Tugas \& Identitas Mahasiswa}]
\small
\begin{tabular}{p{2.2cm}p{5.5cm}p{2.2cm}p{5.0cm}}
\textbf{Nama} & : \dotfill & \textbf{NPM} & : \dotfill \\
\textbf{Prodi/Kelas} & : S1 Teknik Elektro & \textbf{Hari/Tanggal} & : \dotfill \\
\textbf{Dosen} & : Ir. Novalio Daratha, Ph.D. & \textbf{Batas Waktu} & : 1 Minggu \\
\end{tabular}
\end{tcolorbox}

\vspace{0.2cm}

\begin{enumerate}[leftmargin=1.5em, label=\textbf{Soal \arabic*.}]

    \item \textbf{[Level C1 -- Mengingat / \textit{Remembering}]} \hfill \textbf{(Bobot: 10\%)}
    \begin{enumerate}[label=\alph*.]
        \item Tuliskan persamaan Transformasi Fourier Diskrit (\textit{Discrete Fourier Transform} - DFT) $N$-titik!
        \item Apa perbedaan mendasar antara koefisien aproksimasi ($cA$) dan koefisien detail ($cD$) pada dekomposisi Wavelet bertingkat?
    \end{enumerate}

    \vspace{0.3cm}

    \item \textbf{[Level C2 -- Memahami / \textit{Understanding}]} \hfill \textbf{(Bobot: 15\%)}
    \begin{enumerate}[label=\alph*.]
        \item Mengapa Transformasi Fourier standar (FFT) gagal memberikan informasi temporal yang akurat ketika menganalisis gangguan transien hubung singkat (*voltage sag* 3-siklus)?
        \item Jelaskan bagaimana \textit{Heisenberg-Gabor Uncertainty Principle} mendasari kompromi resolusi waktu vs resolusi frekuensi pada STFT dan Wavelet!
    \end{enumerate}

    \vspace{0.3cm}

    \item \textbf{[Level C3 -- Menerapkan / \textit{Applying}]} \hfill \textbf{(Bobot: 20\%)} \\
    Sebuah sinyal arus beban non-linier $i(t)$ disampling dengan frekuensi $f_s = 1000$ Hz sebanyak $N = 1000$ sampel (durasi 1 detik). Spektrum amplitudo FFT $|X[k]|$ menunjukkan puncak dominan pada indeks bin $k = 50$ (amplitudo 10.0 A), $k = 150$ (amplitudo 3.0 A), dan $k = 250$ (amplitudo 1.5 A).
    \begin{enumerate}[label=\alph*.]
        \item Hitunglah frekuensi analog ($f_k$ dalam Hz) yang bersesuaian dengan ketiga bin spektral tersebut!
        \item Tentukan frekuensi fundamental dan orde harmonisa yang muncul.
        \item Hitung nilai \textit{Total Harmonic Distortion} (THD$_i$) dari sinyal arus tersebut:
        \begin{equation*}
            \text{THD}_i = \frac{\sqrt{I_3^2 + I_5^2 + \dots}}{I_1} \times 100\%
        \end{equation*}
    \end{enumerate}

    \newpage

    \item \textbf{[Level C4 -- Menganalisis / \textit{Analyzing}]} \hfill \textbf{(Bobot: 20\%)} \\
    Sinyal getaran bantalan (*bearing*) motor induksi 3-fasa diuraikan menggunakan DWT 3-Level (Wavelet Daubechies \texttt{db4}). Hasil energi pada masing-masing sub-band dihitung sebagai berikut:
    \begin{center}
    \begin{tabular}{|l|c|c|c|}
    \hline
    \rowcolor{boxbg}
    \textbf{Sub-Band} & \textbf{Rentang Frekuensi (Hz)} & \textbf{Kondisi Normal (Energi)} & \textbf{Kondisi Rusak (Energi)} \\ \hline
    Aproksimasi $cA_3$ & $0 - 125$ Hz & $85.4 \text{ J}$ & $84.2 \text{ J}$ \\ \hline
    Detail $cD_3$ & $125 - 250$ Hz & $8.2 \text{ J}$ & $9.1 \text{ J}$ \\ \hline
    Detail $cD_2$ & $250 - 500$ Hz & $4.1 \text{ J}$ & $12.6 \text{ J}$ \\ \hline
    Detail $cD_1$ & $500 - 1000$ Hz & $2.3 \text{ J}$ & $48.5 \text{ J}$ (Lonjakan 21x) \\ \hline
    \end{tabular}
    \end{center}
    \textbf{Tugas Analisis:}
    \begin{enumerate}[label=\alph*.]
        \item Pada sub-band manakah anomali kerusakan mekanis bantalan motor paling awal dan paling jelas terdeteksi?
        \item Mengapa koefisien detail frekuensi tinggi ($cD_1$) lebih sensitif terhadap cacat mikro permukaan bantalan (*spalling/pitting*) dibandingkan koefisien aproksimasi ($cA_3$)?
    \end{enumerate}

    \vspace{0.4cm}

    \item \textbf{[Level C5 -- Mengevaluasi / \textit{Evaluating}]} \hfill \textbf{(Bobot: 15\%)} \\
    Dalam pra-pemrosesan sinyal transien kualitas daya yang terkontaminasi derau putih Gaussian:
    \begin{itemize}
        \item \textbf{Metode A}: Filter Digital Lolos-Rendah Butterworth (IIR) Ordo-4 dengan frekuensi cut-off 200 Hz.
        \item \textbf{Metode B}: \textit{Wavelet Denoising} dengan \textit{Soft-Thresholding} Donoho pada koefisien detail DWT.
    \end{itemize}
    Evaluasilah performa kedua metode dalam mempertahankan ketajaman tepi transien (*edge preserving*) saat terjadi sambaran petir surja hubung!

    \vspace{0.4cm}

    \item \textbf{[Level C6 -- Merancang / \textit{Creating}]} \hfill \textbf{(Bobot: 20\%)} \\
    Rancanglah program Python atau Octave mandiri untuk sistem deteksi gangguan kualitas daya berbasis \textbf{Wavelet Multi-Resolution Analysis}:
    \begin{enumerate}[label=\alph*.]
        \item Membaca sinyal gelombang tegangan 50 Hz berdurasi 10 siklus yang disisipi gangguan *voltage sag* pada siklus ke 4 s.d. 6.
        \item Melakukan dekomposisi DWT hingga level 4 menggunakan basis wavelet \texttt{db4}.
        \item Menghitung nilai ambang deteksi $\tau = \mu_{cD1} + 3\sigma_{cD1}$ untuk mengisolasi titik waktu mula (*start-time*) dan titik waktu akhir (*end-time*) gangguan secara otomatis.
        \item Menampilkan plot 5-panel (Sinyal Asli, $cA_4, cD_3, cD_2, cD_1$) dengan penanda waktu gangguan.
    \end{enumerate}

\end{enumerate}

\end{document}
